% Generated draft: requires mathematical and source review.
\documentclass[11pt]{article}
\usepackage[margin=25mm]{geometry}
\usepackage{fontspec}
\setmainfont{Latin Modern Roman}
\usepackage{amsmath,amssymb,mathtools,unicode-math}
\setmathfont{Latin Modern Math}
\usepackage{enumitem,needspace,xurl,hyperref}
\setlength{\parindent}{0pt}
\setlength{\parskip}{5pt}
\setlength{\emergencystretch}{3em}
\newcommand{\N}{\mathbb{N}}
\newcommand{\Z}{\mathbb{Z}}
\newcommand{\Q}{\mathbb{Q}}
\newcommand{\R}{\mathbb{R}}
\newcommand{\C}{\mathbb{C}}
\newcommand{\F}{\mathbb{F}}
\newcommand{\A}{\mathbb{A}}
\newcommand{\PP}{\mathbb{P}}
\newcommand{\E}{\mathbb{E}}
\newcommand{\eps}{\varepsilon}
\newcommand{\dd}{\,\mathrm{d}}
\begin{document}
\section*{UnsolvedMath ID: 800902}
\textbf{Source problem code:} catalogue entry
\subsection{Definitions and mathematical statement}
For a function \(\varphi\) on a set \(\Omega\), write
\[
\{\varphi<-s\}:=\{z\in\Omega:\varphi(z)<-s\}.
\]
The supplied source describes a preceding estimate as giving optimal exponential integrability, while remaining nonsharp because of a polynomial prefactor. It asks whether that prefactor can be removed: does there exist a finite constant \(c_n\), depending only on \(n\) and possibly different from the constant in the preceding estimate, such that
\[
\operatorname{Vol}\bigl(\{\varphi<-s\}\bigr)
\leq c_n\delta_{\Omega}^{2n}\exp\!\left(-\frac{2ns}{\mu}\right)
\]
holds uniformly for the data and parameter values intended in the source?

\textbf{Unresolved qualifications.} The supplied excerpt omits the preceding estimate and its hypotheses. In particular, it does not specify the meaning or range of \(n\), the ambient space or admissible domains \(\Omega\), the class and normalization of \(\varphi\), the volume measure, the definition of \(\delta_{\Omega}\), the definition and admissible values of \(\mu\), or the range of \(s\). The displayed expression requires \(\mu\neq0\) to be defined, but the source excerpt supplies no hypothesis guaranteeing this. Thus the displayed question preserves the supplied target inequality, but the record does not determine a fully quantified mathematical conjecture.
\subsection{Short English statement}
Question 33 asks whether a polynomial prefactor in a stated exponential sublevel-volume estimate can be removed while retaining the displayed exponential rate. The preceding estimate and the hypotheses defining its admissible data are absent, so the intended uniform inequality requires clarification.
\subsection{Sources}
\begin{itemize}
\item Supplied local UnsolvedMath registry record, permanent ID \(800902\), exact source problem code \texttt{catalogue entry}; title: \emph{Open Problems in Pluripotential Theory --- Question 33 (ACHKZ)}. The registry attributes the problem to Dinew, Guedj, and Zeriahi, with proposed year 2015, and labels it \texttt{open}. This status is supplied registry metadata, not an independently verified mathematical assessment.
\item Dinew, Guedj, and Zeriahi, \emph{Open Problems in Pluripotential Theory} (2015), Question 33. Supplied source URL: \url{https://arxiv.org/abs/1511.00705}. The record reports access on 2026-10-16 and extraction method \texttt{arxiv\_latex}; these are supplied provenance details, not a claim of independent access or verification.
\item The supplied background gives the attribution/context label \texttt{ACHKZ}. Its expansion and any associated bibliographic reference are not included in the record.
\end{itemize}
\end{document}
